\chapter[Wprowadzenie do PDE]{Równania różniczkowe cząstkowe: pojęcia i pierwsze zjawiska}
\label{chap:pde-intro}

Równania różniczkowe cząstkowe (ang. \emph{partial differential equations}, PDE)
opisują związki między pochodnymi funkcji wielu zmiennych. Przy ich pomocy
formułuje się prawa przepływu ciepła, rozchodzenia się fal, ruchu płynów,
dynamiki kwantowej i geometrii rozmaitości. Ten sam operator różniczkowy
może prowadzić do zupełnie różnych zagadnień w zależności od tego, na
jakim obszarze pracujemy, jakie narzucamy warunki i w jakiej klasie
szukamy rozwiązań.

Celem rozdziału jest zbudowanie języka potrzebnego w dalszej części
skryptu. Wprowadzimy rząd i typ nieliniowości równania, symbol główny,
powierzchnie charakterystyczne i podstawowe klasy równań. Następnie,
na przykładach, oddzielimy trzy problemy: istnienie rozwiązania,
jego jednoznaczność i stabilną zależność od danych.

\section{Co jest niewiadomą w równaniu PDE?}
\label{sec:pde-unknown}

Niech $n\geq1$, $\Omega\subset\R^n$ będzie niepustym zbiorem otwartym,
a $u\colon\Omega\to\R$ funkcją niewiadomą. Równanie różniczkowe
cząstkowe wiąże wartości $u$ i jej pochodnych cząstkowych. Przykładowo:
\begin{align}
 -\Delta u&=f &&\text{(równanie Poissona)},\label{eq:poisson-intro}\\
 u_t-\kappa\Delta_xu&=0,\quad \kappa>0
 &&\text{(równanie ciepła)},\label{eq:heat-intro}\\
 u_{tt}-c^2\Delta_xu&=0,\quad c>0
 &&\text{(równanie falowe)},\label{eq:wave-intro}\\
 u_t+u\,u_x&=0 &&\text{(równanie Burgersa bez lepkości)}.
 \label{eq:burgers-intro}
\end{align}
W równaniach ciepła i falowym niewiadoma zależy od czasu $t$ i
zmiennych przestrzennych $x=(x_1,\ldots,x_d)$; w podanym równaniu
Burgersa jest jedna zmienna przestrzenna, czyli $d=1$. Zapis $\Delta_x$ oznacza
laplasjan liczony wyłącznie względem zmiennych przestrzennych, tj.
$\Delta_xu=\sum_{j=1}^d\partial_{x_j}^2u$.

\begin{uwaga}[Równanie a zagadnienie z dodatkowymi warunkami]
Jeśli na niepustym zbiorze otwartym istnieje rozwiązanie klasyczne
$u$ równania $-\Delta u=f$, to bez dodatkowych warunków nie jest
ono jedyne: dla każdej stałej $b\in\R$ funkcja $u+b$ również
spełnia to równanie. Dopiero po
określeniu obszaru, warunków brzegowych i klasy funkcji otrzymujemy
konkretne \emph{zagadnienie}. Przykładowo każda funkcja afiniczna
$u(x)=a\cdot x+b$ jest harmoniczna w $\R^n$; nie można zatem mówić
o jednoznaczności bez dalszych założeń.
\end{uwaga}

W pierwszych rozdziałach przez rozwiązanie klasyczne równania rzędu
$m$ rozumiemy funkcję klasy $C^m$ spełniającą równanie punktowo. Dla równań
ewolucyjnych wygodniej jest osobno liczyć pochodne czasowe i
przestrzenne; na przykład $C^{1,2}$ dla kolejności argumentów $(t,x)$ oznacza ciągłość funkcji, jej pierwszej pochodnej czasowej
oraz wszystkich pochodnych przestrzennych rzędu pierwszego i drugiego.
Nie wymaga się tu istnienia pochodnych mieszanych czasowo-przestrzennych. Rozwiązania słabe i dystrybucyjne
wprowadzimy później.

\section{Rząd równania i rodzaje nieliniowości}
\label{sec:pde-types}

Niech $\N_0=\{0,1,2,\ldots\}$ i $n\geq1$.
Dla wielowskaźnika $\alpha=(\alpha_1,\ldots,\alpha_n)\in\N_0^n$
przyjmujemy $|\alpha|=\alpha_1+\cdots+\alpha_n$ i
$\partial^\alpha=\partial_{x_1}^{\alpha_1}\cdots
\partial_{x_n}^{\alpha_n}$. Niech $m\geq1$, a $J^m u(x)$ oznacza
uporządkowaną rodzinę wszystkich wartości $\partial^\alpha u(x)$,
$|\alpha|\leq m$ (tzw. $m$-dżet funkcji w punkcie).

\begin{definicja}[Równanie rzędu $m$]
Niech $n,m\geq1$ będą liczbami całkowitymi,
$\Omega\subset\R^n$ będzie otwarta, a
$N_m=\#\{\alpha\in\N_0^n:|\alpha|\leq m\}$.
Niech $U\subset\Omega\times\R^{N_m}$ będzie zbiorem otwartym
i $F\colon U\to\R$. Dla funkcji $u\in C^m(\Omega,\R)$
takiej, że $(x,J^m u(x))\in U$ dla każdego $x\in\Omega$,
równanie $F(x,J^m u(x))=0$ nazywamy równaniem rzędu $m$,
jeżeli $F$ zależy nietrywialnie od współrzędnych $z_\alpha$
o $|\alpha|=m$, tj. istnieją dwa punkty zbioru $U$ o tym samym
$x$ i tych samych współrzędnych $z_\beta$ dla $|\beta|<m$,
w których $F$ przyjmuje różne wartości. Pojęcie rzędu odnosi się do przedstawionego równania.
W szczególnych punktach część najwyższego rzędu może znikać;
na przykład operator $x\,\partial_x^2+\partial_x$ ma rząd $2$
na $\R$, choć w punkcie $x=0$ jego symbol drugiego rzędu
zeruje się dla każdego kowektora.
\end{definicja}

\begin{definicja}[Liniowość, półliniowość i quasiliniowość]
Niech $\Omega\subset\R^n$ będzie otwarta i $m\geq1$.
\begin{enumerate}
\item Równanie \emph{liniowe} rzędu $m$ ma postać
$Lu=f$, gdzie $Lu=\sum_{|\alpha|\leq m}a_\alpha(x)\partial^\alpha u$,
a współczynniki $a_\alpha$ oraz $f$ nie zależą od $u$.
Jest \emph{jednorodne}, jeśli $f=0$.
\item Równanie jest \emph{półliniowe} (semiliniowe), jeśli ma postać
$\sum_{|\alpha|=m}a_\alpha(x)\partial^\alpha u
+G(x,J^{m-1}u)=0$. Pochodne najwyższego rzędu wchodzą zatem
liniowo, a ich współczynniki zależą tylko od $x$.
\item Równanie jest \emph{quasiliniowe}, jeśli można je zapisać jako
$\sum_{|\alpha|=m}a_\alpha(x,J^{m-1}u)\partial^\alpha u
+G(x,J^{m-1}u)=0$.
\item Równanie nazywamy \emph{w pełni nieliniowym}, gdy jego zależność
od pochodnych najwyższego rzędu nie jest afiniczna.
\end{enumerate}
Terminy te dotyczą przyjętego rzędu i postaci równania; równania
liniowe są szczególnym przypadkiem półliniowych, a półliniowe
szczególnym przypadkiem quasiliniowych. Klasyfikacja odnosi się do
\emph{wyrażenia różniczkowego}, nie wyłącznie do zbioru jego
rozwiązań: równania $\Delta u=0$ i $(\Delta u)^2=0$ mają ten
sam zbiór rozwiązań klasycznych, ale inne postaci algebraiczne.
\end{definicja}

\begin{przyklad}[Rozpoznawanie typu]
Dla niewiadomej dwóch zmiennych, oznaczanych przez $(x,y)$
w pierwszych czterech przykładach, a przez $(t,x)$ w ostatnim:
\begin{enumerate}
\item $u_{xx}+u_{yy}=\sin x$ jest liniowe i niejednorodne.
\item $u_{xx}+u_{yy}+u^3=0$ jest półliniowe.
\item $(1+u_x^2)u_{xx}+u_{yy}=0$ jest quasiliniowe: współczynnik
przy $u_{xx}$ zależy od pochodnej pierwszego rzędu.
\item $u_{xx}u_{yy}-u_{xy}^2=1$ jest w pełni nieliniowe
(równanie typu Monge'a--Ampère'a).
\item $u_t+u\,u_x=0$ jest quasiliniowe pierwszego rzędu.
\end{enumerate}
\end{przyklad}

\begin{stwierdzenie}[Zasada superpozycji]
Niech $V,W$ będą przestrzeniami wektorowymi funkcji nad $\R$,
$L\colon V\to W$ operatorem liniowym, $u_1,u_2\in V$
oraz $f_1,f_2,f\in W$. Jeśli $Lu_1=f_1$ i
$Lu_2=f_2$, to dla dowolnych $a,b\in\R$ zachodzi
$L(au_1+bu_2)=af_1+bf_2$. W szczególności zbiór rozwiązań
równania jednorodnego $Lu=0$ jest podprzestrzenią wektorową
$V$, a zbiór rozwiązań $Lu=f$, jeżeli jest niepusty,
ma postać $u_0+\ker L$.
\end{stwierdzenie}
\begin{proof}
Z liniowości $L$ otrzymujemy $L(au_1+bu_2)=aLu_1+bLu_2
=af_1+bf_2$. Zbiór $\ker L$ zawiera wektor zerowy i jest zamknięty
na kombinacje liniowe, więc jest podprzestrzenią wektorową $V$.
Jeżeli $Lu_0=f$, to dla $u\in V$ mamy $Lu=f$ wtedy i tylko wtedy,
gdy $L(u-u_0)=0$, czyli $u-u_0\in\ker L$. Zbiór rozwiązań
jest zatem dokładnie równy $u_0+\ker L$.
\end{proof}

\begin{uwaga}
Dla równania nieliniowego nie wolno zazwyczaj dodawać rozwiązań.
Rozważmy na przykład równanie $w_x=w^2$ na
$\Omega=(-1,1)^2$, gdzie $w=w(x,y)$.
Funkcje $w_1(x,y)=1/(2-x)$ i $w_2(x,y)=1/(3-x)$
są rozwiązaniami, ponieważ $w_j'=w_j^2$.
Ich suma nie jest jednak rozwiązaniem, gdyż
$(w_1+w_2)'=w_1^2+w_2^2\ne(w_1+w_2)^2$.
Podobnie $2w_1$ nie jest rozwiązaniem: pokazuje to,
że nie działa również zwykłe mnożenie rozwiązań przez skalar.
\end{uwaga}

\section{Dane początkowe i brzegowe}
\label{sec:pde-data}

Niech $\Omega\subset\R^d$ będzie obszarem ograniczonym, tj.
niepustym, spójnym zbiorem otwartym, i załóżmy, że jej brzeg
$\partial\Omega$ jest co najmniej klasy $C^1$ tam, gdzie używamy
wektora normalnego $\nu$ skierowanego na zewnątrz.

\begin{definicja}[Podstawowe warunki brzegowe]
Niech $\Omega\subset\R^d$ będzie obszarem ograniczonym o brzegu
klasy $C^1$, $g,\beta\colon\partial\Omega\to\R$ będą zadanymi
funkcjami, a $\nu$ jednostkowym wektorem normalnym skierowanym
na zewnątrz. Dla rzeczywistej funkcji $u$ ciągłej na
$\overline\Omega$ (w warunkach Neumanna i Robina dodatkowo
klasy $C^1$ aż do brzegu) rozważamy:
\begin{enumerate}
\item \emph{warunek Dirichleta}: $u|_{\partial\Omega}=g$;
\item \emph{warunek Neumanna}: $\partial_\nu u|_{\partial\Omega}=g$,
 gdzie $\partial_\nu u=\nabla u\cdot\nu$;
\item \emph{warunek Robina}: $\partial_\nu u+\beta u=g$
 na $\partial\Omega$, gdzie $\beta$ jest zadaną funkcją.
\end{enumerate}
Warunki te rozumiemy na razie punktowo, dla funkcji o odpowiedniej
regularności aż do brzegu.
\end{definicja}

Dla zagadnień ewolucyjnych na $(0,T)\times\Omega$, gdzie $T>0$,
warunek początkowy zapisujemy na przykład jako $u(0,x)=u_0(x)$.
Równanie falowe drugiego rzędu względem czasu wymaga zwykle także
warunku $u_t(0,x)=u_1(x)$. Jeśli $\Omega$ jest ograniczona,
potrzebne są ponadto odpowiednie warunki na bocznym brzegu
$(0,T)\times\partial\Omega$.

\begin{przyklad}[Zgodność danych na krawędzi czasoprzestrzeni]
Niech $T>0$, $\Omega\subset\R^d$ będzie obszarem ograniczonym,
$u_0\colon\overline\Omega\to\R$ oraz
$g\colon[0,T]\times\partial\Omega\to\R$.
Jeżeli $u\in C([0,T]\times\overline\Omega)$ spełnia
$u(0,x)=u_0(x)$ dla $x\in\overline\Omega$ oraz
$u(t,x)=g(t,x)$ dla $(t,x)\in[0,T]\times\partial\Omega$,
to koniecznie $u_0(x)=g(0,x)$ dla każdego $x\in\partial\Omega$.
Jest to najprostszy \emph{warunek zgodności} danych; wyższa
regularność często narzuca dalsze warunki.
\end{przyklad}

\begin{uwaga}[Zagadnienie Cauchy'ego]
Przez zagadnienie Cauchy'ego rozumiemy poszukiwanie rozwiązania
na podstawie danych zadanych na odpowiedniej hiperpłaszczyźnie
lub hiperpowierzchni. Nie należy utożsamiać go z zagadnieniem
Dirichleta. Na przykład zadanie wartości $u$ i pochodnej normalnej
na części brzegu dla równania Laplace'a jest zagadnieniem Cauchy'ego,
które, jak zobaczymy w~\S\ref{sec:pde-hadamard}, jest niestabilne.
\end{uwaga}

\section{Symbol główny i powierzchnie charakterystyczne}
\label{sec:pde-symbol}

Pojęcie symbolu ujawnia, co z równania pozostaje przy badaniu
oscylacji o bardzo dużej częstości. Niech $m\geq1$,
$\Omega\subset\R^n$ będzie otwarta, a
$P=\sum_{|\alpha|\leq m}a_\alpha(x)\partial^\alpha$
liniowym operatorem różniczkowym o współczynnikach ciągłych
na $\Omega$. Zakładamy, że przynajmniej jeden współczynnik
najwyższego rzędu nie znika identycznie.

\begin{definicja}[Symbol główny]
\emph{Symbolem głównym} operatora
$P=\sum_{|\alpha|\leq m}a_\alpha(x)\partial^\alpha$ rzędu $m$
na otwartym zbiorze $\Omega\subset\R^n$ nazywamy funkcję
$p_m\colon\Omega\times\R^n\to\C$, będącą dla każdego ustalonego $x$
wielomianem jednorodnym stopnia $m$ względem $\xi$ (dopuszczamy
wielomian zerowy w punktach zaniku części głównej), daną wzorem
\begin{equation}\label{eq:principal-symbol}
 p_m(x,\xi)=\sum_{|\alpha|=m}a_\alpha(x)\xi^\alpha.
\end{equation}
W naszej konwencji $\partial^\alpha$ nie zawiera czynników $-i$.
Jeśli ten sam operator zapiszemy w konwencji $D_j=-i\partial_{x_j}$,
to $\partial^\alpha=i^{|\alpha|}D^\alpha$, więc jego symbol główny
w konwencji $D$ jest równy $i^mp_m(x,\xi)$.
\end{definicja}

\begin{stwierdzenie}[Test wysokich częstości]
Niech $\Omega\subset\R^n$ będzie otwarta, $m\geq1$ oraz
$P=\sum_{|\alpha|\leq m}a_\alpha(x)\partial^\alpha$ będzie liniowym
operatorem rzędu $m$ o ciągłych współczynnikach, działającym na
funkcjach zespolonych (w razie potrzeby po zespoleniu operatora
rzeczywistego). Niech $p_m$ będzie jego symbolem głównym
w konwencji~\eqref{eq:principal-symbol}, a $x\in\Omega$, $\xi\in\R^n$. Dla $\lambda>0$ oraz
$v_\lambda(z)=e^{i\lambda z\cdot\xi}$ zachodzi
\[
 e^{-i\lambda x\cdot\xi}Pv_\lambda(x)
 =i^m\lambda^m p_m(x,\xi)+O(\lambda^{m-1})
 \quad\text{gdy }\lambda\to\infty.
\]
Stała w symbolu $O(\lambda^{m-1})$ może zależeć od $x$ i $\xi$.
\end{stwierdzenie}
\begin{proof}
Mamy $\partial^\alpha v_\lambda(x)=
(i\lambda)^{|\alpha|}\xi^\alpha v_\lambda(x)$.
Po podstawieniu do definicji $P$ wyrazy z $|\alpha|=m$
dają dokładnie $i^m\lambda^mp_m(x,\xi)$, a skończona
suma pozostałych wyrazów jest rzędu co najwyżej $\lambda^{m-1}$.
\end{proof}

\begin{definicja}[Eliptyczność]
Liniowy operator skalarny $P$ rzędu $m$ na zbiorze otwartym
$\Omega\subset\R^n$, o symbolu głównym $p_m$ danym przez
\eqref{eq:principal-symbol}, nazywamy \emph{eliptycznym w punkcie}
$x\in\Omega$, jeżeli $p_m(x,\xi)\ne0$ dla każdego
$\xi\in\R^n\setminus\{0\}$. Jest \emph{eliptyczny na $\Omega$},
gdy warunek zachodzi w każdym punkcie.
\end{definicja}

\begin{przyklad}
Dla $P=-\Delta=-\sum_{j=1}^n\partial_{x_j}^2$
przyjęta konwencja daje $p_2(x,\xi)=-|\xi|^2$.
Symbol nie zeruje się poza $\xi=0$, więc laplasjan jest
operatorem eliptycznym. Znak symbolu nie ma znaczenia
dla samej definicji eliptyczności.
\end{przyklad}

\begin{uwaga}[Eliptyczność a jednostajna eliptyczność]
Dla operatora drugiego rzędu o części głównej
$-\sum_{i,j=1}^n a_{ij}(x)\partial_{x_ix_j}$, gdzie
$A(x)=(a_{ij}(x))_{i,j=1}^n$ jest macierzą rzeczywistą i symetryczną,
warunek jednostajnej eliptyczności (w konwencji dolnego
ograniczenia) ma postać
$\xi^{\mathsf T}A(x)\xi\ge\lambda|\xi|^2$
dla wszystkich $x\in\Omega$, $\xi\in\R^n$ i pewnej
\emph{jednej} stałej $\lambda>0$. Samo niezerowanie się
symbolu głównego w poszczególnych punktach nie daje
takiej stałej. Przykładowo operator
$P=-x^2\partial_x^2$ jest eliptyczny na $(0,1)$,
ale nie jest jednostajnie eliptyczny: dla $x\downarrow0$
wartość współczynnika $x^2$ dąży do zera.
Wymóg dolnego oszacowania przez $\lambda|\xi|^2$
dotyczy operatorów drugiego rzędu w postaci dodatnio
określonej; nie jest definicją wszystkich eliptycznych
operatorów dowolnego rzędu.
\end{uwaga}

\begin{definicja}[Powierzchnia charakterystyczna]
Niech $P$ będzie liniowym operatorem skalarnym rzędu $m$ na
$\Omega\subset\R^n$, a $\Sigma\subset\Omega$ hiperpowierzchnią klasy $C^1$.
W otoczeniu otwartym $U\subset\Omega$ rozważanego punktu przyjmijmy
$\Sigma\cap U=\{z\in U:\phi(z)=0\}$, gdzie
$\phi\in C^1(U)$ i $\mathrm d\phi(z)\ne0$ dla $z\in\Sigma\cap U$.
Mówimy, że $\Sigma$ jest \emph{charakterystyczna dla $P$
w punkcie} $x\in\Sigma$, jeżeli
$p_m(x,\mathrm d\phi(x))=0$. W przeciwnym razie jest
\emph{niecharakterystyczna w $x$}. Warunek nie zależy od
wyboru lokalnej funkcji definiującej $\phi$.
\end{definicja}

\begin{proof}[Uzasadnienie niezależności od funkcji definiującej]
Jeśli $\phi$ i $\psi$ są dwiema lokalnymi funkcjami definiującymi
tę samą hiperpowierzchnię klasy $C^1$, to
$T_x\Sigma=\ker\mathrm d\phi(x)=\ker\mathrm d\psi(x)$.
Istotnie, różniczkowanie funkcji definiującej wzdłuż krzywej
na $\Sigma$ daje zero, a zarówno przestrzeń styczna, jak i jądro
niezerowej różniczki mają wymiar $n-1$. Wybierzmy $v\in\R^n$
tak, by $\mathrm d\phi(x)v=1$. Dla dowolnego $w\in\R^n$
wektor $w-\mathrm d\phi(x)w\,v$ należy do wspólnego jądra,
więc $\mathrm d\psi(x)w=\mathrm d\psi(x)v\,\mathrm d\phi(x)w$.
Stąd $\mathrm d\psi(x)=a\,\mathrm d\phi(x)$, gdzie
$a=\mathrm d\psi(x)v\ne0$, ponieważ $\mathrm d\psi(x)\ne0$.
Ponieważ $p_m(x,\cdot)$ jest jednorodny stopnia $m$,
$p_m(x,\mathrm d\psi(x))=a^mp_m(x,\mathrm d\phi(x))$.
Oba wyrażenia zerują się zatem jednocześnie.
\end{proof}

\begin{przyklad}[Stożek charakterystyczny równania falowego]
Dla operatora $P=\partial_t^2-c^2\Delta_x$ w $\R^{1+d}$
symbol główny jest równy
$p_2(t,x;\tau,\xi)=\tau^2-c^2|\xi|^2$.
Kowektory $(\tau,\xi)\ne0$ spełniające
$\tau^2=c^2|\xi|^2$ są charakterystyczne. W przypadku
jednego wymiaru przestrzennego proste $t-x/c=\mathrm{const}$
i $t+x/c=\mathrm{const}$ są charakterystyczne.
Natomiast powierzchnia początkowa $t=0$ jest
niecharakterystyczna, bo $p_2(1,0)=1$.
\end{przyklad}

\begin{uwaga}[Pułapka: symbol równania ciepła]
Dla $P=\partial_t-\kappa\Delta_x$, traktowanego jako zwykły
operator drugiego rzędu we wszystkich zmiennych $(t,x)$,
symbol główny wynosi $-\kappa|\xi|^2$: nie zależy od $\tau$.
W takim sensie powierzchnia $t=0$ jest charakterystyczna
($p_2((t,x),(1,0))=0$),
chociaż zadanie początkowe dla równania ciepła jest naturalne
i stabilne w odpowiednich przestrzeniach. Dlatego równania
paraboliczne klasyfikuje się przy użyciu \emph{ważonego}
porządku pochodnych: jednej pochodnej czasowej przypisuje się
wagę $2$, a przestrzennej wagę $1$. Przy konwencji symbolu
$D_j=-i\partial_{x_j}$ prowadzi to dla operatora ciepła do
symbolu parabolicznego $i\tau+\kappa|\xi|^2$, który dla
rzeczywistych $(\tau,\xi)\ne0$ nie znika. W szczególności
nie wolno bez dodatkowych zastrzeżeń stosować do równania
ciepła intuicji charakterystyk równania falowego.
\end{uwaga}

\section{Klasyfikacja operatorów drugiego rzędu}
\label{sec:pde-classification}

Dla równania drugiego rzędu w dwóch zmiennych rzeczywistych
rozważmy część główną
\begin{equation}\label{eq:2d-principal}
 A(x,y)u_{xx}+2B(x,y)u_{xy}+C(x,y)u_{yy}.
\end{equation}
Współczynniki $A,B,C$ są rzeczywiste i ciągłe, a odpowiadająca
im macierz symetryczna to
$\mathsf A(x,y)=\begin{pmatrix}A&B\\B&C\end{pmatrix}$.

\begin{definicja}[Klasyfikacja w dwóch zmiennych]
Niech $\Omega\subset\R^2$ będzie otwarta, $A,B,C\colon\Omega\to\R$
będą ciągłe, a $\mathsf A=\begin{pmatrix}A&B\\B&C\end{pmatrix}$
będzie macierzą części głównej~\eqref{eq:2d-principal}.
Ustalmy punkt $(x,y)\in\Omega$.
Oznaczmy $\delta(x,y)=B(x,y)^2-A(x,y)C(x,y)$.
Równanie z częścią główną~\eqref{eq:2d-principal} ma w tym punkcie
typ \emph{eliptyczny}, gdy $\delta<0$;
\emph{hiperboliczny}, gdy $\delta>0$;
\emph{paraboliczny w tradycyjnej klasyfikacji dwuzmiennowej},
gdy $\delta=0$ i $\mathsf A\ne0$.
Jeśli $\mathsf A=0$, część drugiego rzędu zanika i tej
klasyfikacji tam nie stosujemy.
\end{definicja}

\begin{stwierdzenie}[Interpretacja geometryczna i niezmienniczość typu]
Niech $\mathsf A=\begin{pmatrix}A&B\\B&C\end{pmatrix}$
będzie rzeczywistą, symetryczną macierzą $2\times2$, różną od zera.
Wtedy:
\begin{enumerate}
\item $B^2-AC<0$ wtedy i tylko wtedy, gdy $\mathsf A$ jest
 dodatnio albo ujemnie określona;
\item $B^2-AC>0$ wtedy i tylko wtedy, gdy $\mathsf A$ ma
 wartości własne przeciwnych znaków;
\item $B^2-AC=0$ wtedy i tylko wtedy, gdy $\mathsf A$ ma rząd $1$.
\end{enumerate}
Pod lokalną zmianą współrzędnych klasy $C^2$ o nieosobliwym
jakobianie liczba dodatnich, ujemnych i zerowych wartości własnych
macierzy części głównej nie zmienia się. W konsekwencji typ
równania jest niezmienniczy względem takich zmian współrzędnych.
\end{stwierdzenie}
\begin{proof}
Macierz $\mathsf A$ ma rzeczywiste wartości własne
$\lambda_1,\lambda_2$, przy czym
$\lambda_1\lambda_2=\det\mathsf A=AC-B^2=-\delta$.
Jeśli $\delta<0$, iloczyn jest dodatni i obie niezerowe
wartości własne mają ten sam znak. Jeśli $\delta>0$,
mają znaki przeciwne. Gdy $\delta=0$, jedna wartość własna
jest zerowa, a druga niezerowa, ponieważ $\mathsf A\ne0$.

Niech $z=\Phi(x)$ będzie lokalną zmianą współrzędnych,
$J=D\Phi(x)$ jej nieosobliwym jakobianem i
$u(x)=v(\Phi(x))$, gdzie $\Phi$ i $v$ są klasy $C^2$.
Piszemy $J_{ki}=\partial_{x_i}\Phi_k(x)$.
Dwukrotne zastosowanie reguły łańcuchowej daje
\[
 \partial_{x_i x_j}u(x)=
 \sum_{k,l=1}^2(\partial_{z_k z_l}v)(\Phi(x))J_{ki}J_{lj}
 +\sum_{k=1}^2(\partial_{z_k}v)(\Phi(x))
       \partial_{x_i x_j}\Phi_k(x).
\]
Jeśli $a_{ij}$ są elementami $\mathsf A$, to w zapisie
$\sum_{k,l=1}^2 b_{kl}\partial_{z_k z_l}v$ otrzymujemy
$b_{kl}=\sum_{i,j=1}^2J_{ki}a_{ij}J_{lj}=(J\mathsf A J^{\mathsf T})_{kl}$.
Macierz $(b_{kl})$ jest symetryczna; po zebraniu pochodnych
mieszanych $\partial_{z_1z_2}v=\partial_{z_2z_1}v$ ich łączny
współczynnik wynosi $2b_{12}$.
Drugi składnik wzoru zawiera jedynie pierwsze pochodne $v$.
Prawo bezwładności
Sylvestera zapewnia równość liczby dodatnich, ujemnych
i zerowych wartości własnych obu macierzy.
\end{proof}

\begin{przyklad}[Trzy podstawowe typy]
\begin{enumerate}
\item $u_{xx}+u_{yy}=0$: $A=C=1$, $B=0$, więc $\delta=-1$;
 typ eliptyczny.
\item $u_{xx}-u_{yy}=0$: $A=1$, $C=-1$, $B=0$, więc
 $\delta=1$; typ hiperboliczny.
\item $u_{xx}+2u_{xy}+u_{yy}=0$: $A=B=C=1$, więc
 $\delta=0$; typ paraboliczny w klasyfikacji dwuzmiennowej.
 Część główna ma postać $(\partial_x+\partial_y)^2$.
\end{enumerate}
\end{przyklad}

\begin{przyklad}[Równanie zmieniające typ]
Rozważmy na $\R^2$ równanie Tricomiego
\begin{equation}\label{eq:tricomi}
 u_{xx}+x\,u_{yy}=0.
\end{equation}
Mamy $A=1$, $B=0$, $C=x$, zatem $\delta=-x$. Równanie jest
eliptyczne w obszarze $x>0$, hiperboliczne w obszarze $x<0$
i typu parabolicznego na prostej $x=0$.
\end{przyklad}

\begin{center}
\begin{tikzpicture}[scale=.9,>=Stealth]
 \fill[gray!10] (-4,-1.25) rectangle (0,1.25);
 \fill[gray!23] (0,-1.25) rectangle (4,1.25);
 \draw[->] (-4.3,0)--(4.45,0) node[right] {$x$};
 \draw[->] (0,-1.5)--(0,1.5);
 \draw[line width=1.1pt] (0,-1.25)--(0,1.25);
 \node[align=center] at (-2,0.65) {hiperboliczny\\$x<0$};
 \node[align=center] at (2,0.65) {eliptyczny\\$x>0$};
 \node[above,align=center] at (0,1.6) {$x=0$: typ paraboliczny};
\end{tikzpicture}
\par\smallskip\small
Typ równania Tricomiego zmienia się na prostej $x=0$.
\end{center}

\begin{uwaga}[Dwa znaczenia słowa „paraboliczny”]
Typ paraboliczny w klasyfikacji wyrażenia
$Au_{xx}+2Bu_{xy}+Cu_{yy}$ dotyczy macierzy części
drugiego rzędu w \emph{dwóch równorzędnych zmiennych}.
Równanie ciepła $u_t-\kappa\Delta_xu=0$ jest natomiast
parabolicznym \emph{równaniem ewolucyjnym} z wyróżnionym
czasem i z innym skalowaniem pochodnych. Terminy są
powiązane historycznie, ale nie należy utożsamiać tych
definicji bez wskazania konwencji.
\end{uwaga}

W większej liczbie zmiennych sama nieokreśloność macierzy
części głównej nie wystarcza do pełnej definicji
hiperboliczności. Dla równania falowego wyróżnia się
kowektor czasowy oraz wymaga odpowiednich własności
rzeczywistych pierwiastków wielomianu symbolu w kierunku
czasowym. Systematyczną teorię rozwinie rozdział o
równaniach hiperbolicznych.

\section{Pierwsze modele: co można udowodnić od razu?}
\label{sec:pde-models}

\subsection{Laplasjan: zasada maksimum i jednoznaczność}

\begin{twierdzenie}[Słaba zasada maksimum]
\label{thm:laplace-max}
Niech $\Omega\subset\R^n$ będzie niepustym, ograniczonym
zbiorem otwartym i niech $u\in C^2(\Omega)\cap
C(\overline\Omega)$ będzie funkcją rzeczywistą. Jeśli
$\Delta u\geq0$ w $\Omega$, to
\begin{equation}\label{eq:maximum-principle}
 \max_{\overline\Omega}u\leq\max_{\partial\Omega}u.
\end{equation}
Jeśli $\Delta u\leq0$, to
$\min_{\overline\Omega}u\geq\min_{\partial\Omega}u$.
W szczególności, gdy $\Delta u=0$, oba wnioski zachodzą.
\end{twierdzenie}
\begin{proof}
Dla $\eps>0$ określmy $u_\eps(x)=u(x)+\eps|x|^2$.
W $\Omega$ zachodzi
$\Delta u_\eps=\Delta u+2n\eps>0$.
Funkcja $u_\eps$ jest ciągła na zwartym zbiorze
$\overline\Omega$, zatem osiąga maksimum. Gdyby osiągała
je w punkcie wewnętrznym, jej hesjan byłby w tym punkcie
ujemnie półokreślony, stąd $\Delta u_\eps\leq0$,
co daje sprzeczność. Maksimum $u_\eps$ leży zatem na brzegu.
Wobec tego dla każdego $x\in\overline\Omega$ mamy
$u(x)+\eps|x|^2\leq
\max_{z\in\partial\Omega}(u(z)+\eps|z|^2)$.
Niech $M=\max_{z\in\overline\Omega}|z|^2<\infty$.
Prawa strona jest nie większa niż
$\max_{\partial\Omega}u+\eps M$, a $\eps|x|^2\geq0$.
Stąd $u(x)\leq\max_{\partial\Omega}u+\eps M$ dla każdego
$\eps>0$. Przechodząc z $\eps$ do zera, otrzymujemy
$u(x)\leq\max_{\partial\Omega}u$; biorąc maksimum po $x$,
dostajemy~\eqref{eq:maximum-principle}.
Dla $\Delta u\leq0$ stosujemy wykazany wynik do $-u$.
\end{proof}

\begin{wniosek}[Jednoznaczność i stabilność zagadnienia Dirichleta]
\label{cor:dirichlet-stability}
Niech $\Omega\subset\R^n$ będzie niepustym ograniczonym
zbiorem otwartym, a $u,v\in C^2(\Omega)\cap
C(\overline\Omega)$ będą funkcjami rzeczywistymi spełniającymi $\Delta u=\Delta v$
w $\Omega$. Wówczas
\[
 \max_{x\in\overline\Omega}|u(x)-v(x)|
 \leq\max_{x\in\partial\Omega}|u(x)-v(x)|.
\]
W szczególności dwa rozwiązania klasyczne zagadnienia
$-\Delta u=f$ z takimi samymi wartościami brzegowymi
muszą być równe.
\end{wniosek}
\begin{proof}
Różnica $w=u-v$ jest harmoniczna. Twierdzenie
\ref{thm:laplace-max}, zastosowane do $w$ i $-w$,
daje dla każdego $x\in\overline\Omega$ nierówności
$w(x)\leq\max_{\partial\Omega}w$ oraz
$-w(x)\leq\max_{\partial\Omega}(-w)$.
Obie prawe strony są nie większe niż
$M=\max_{\partial\Omega}|w|$, więc $|w(x)|\leq M$.
Jeśli wartości brzegowe $u$ i $v$ są równe, to $M=0$
i $u=v$ na $\overline\Omega$.
\end{proof}

Zauważmy, że udowodniliśmy \emph{jednoznaczność i stabilność
względem wartości brzegowych przy ustalonej prawej stronie $f$},
ale nie istnienie rozwiązania Dirichleta dla dowolnego obszaru
i dowolnych danych. Oszacowanie nie dotyczy zmiany prawej
strony równania. Istnienie wymaga osobnej argumentacji.

\subsection{Warunek Neumanna: brak jednoznaczności i warunek zgodności}

\begin{stwierdzenie}[Konieczny warunek zgodności]
Niech $n\geq1$, $\Omega\subset\R^n$ będzie obszarem ograniczonym
o brzegu klasy $C^1$, a $f\in C(\overline\Omega,\C)$ i
$g\in C(\partial\Omega,\C)$. Niech $\nu$ będzie jednostkową
normalną zewnętrzną. Jeśli $u\in C^2(\overline\Omega,\C)$ spełnia
$-\Delta u=f$ w $\Omega$ oraz $\partial_\nu u=g$
na $\partial\Omega$, to
\begin{equation}\label{eq:neumann-compatibility}
 \int_\Omega f(x)\,\mathrm d\lambda_n(x)
 +\int_{\partial\Omega}g(x)\,\mathrm d\sigma(x)=0.
\end{equation}
Tutaj $\lambda_n$ oznacza miarę Lebesgue'a na $\R^n$,
a $\sigma$ miarę powierzchniową na $\partial\Omega$.
Jeśli $u,v\in C^2(\overline\Omega,\C)$ rozwiązują to samo
zagadnienie Neumanna, to funkcja $u-v$ jest stała na $\Omega$.
\end{stwierdzenie}
\begin{proof}
Ponieważ $\nabla u$ jest polem klasy $C^1$ aż do brzegu,
a $\Omega$ jest obszarem ograniczonym o brzegu klasy $C^1$,
możemy zastosować twierdzenie Gaussa do pola $\nabla u$.
Daje ono
$\int_\Omega\Delta u(x)\,\mathrm d\lambda_n(x)
=\int_{\partial\Omega}\partial_\nu u(x)\,\mathrm d\sigma(x)$.
Po podstawieniu $\Delta u=-f$ oraz $\partial_\nu u=g$
i przeniesieniu obu całek na jedną stronę otrzymujemy
\eqref{eq:neumann-compatibility}.

Dla $w=u-v$ zachodzi $\Delta w=0$ i $\partial_\nu w=0$.
Stosując twierdzenie Gaussa do pola $\overline w\nabla w$ klasy $C^1$
(w przypadku zespolonym osobno do części rzeczywistej i urojonej)
i korzystając z równości
$\operatorname{div}(\overline w\nabla w)=|\nabla w|^2+\overline w\Delta w$,
otrzymujemy pierwszą tożsamość Greena:
\[
 \int_\Omega|\nabla w(x)|^2\,\mathrm d\lambda_n(x)
 =\int_{\partial\Omega}\overline{w(x)}\partial_\nu w(x)\,
 \mathrm d\sigma(x)
 -\int_\Omega\overline{w(x)}\Delta w(x)\,\mathrm d\lambda_n(x)=0.
\]
Gdyby $|\nabla w(x_0)|>0$ dla pewnego $x_0\in\Omega$,
to z ciągłości istniałaby kula zawarta w $\Omega$, na której
$|\nabla w|^2$ byłaby ograniczona z dołu przez dodatnią stałą.
Całka byłaby wtedy dodatnia, wbrew otrzymanej równości.
Zatem $\nabla w=0$ w $\Omega$. Na każdej kuli zawartej
w $\Omega$ funkcja $w$ jest stała, co wynika z całkowania
pochodnej wzdłuż odcinków. Jest więc lokalnie stała,
a spójność $\Omega$ implikuje jej stałość na całym obszarze.
Odwrotnie, dodanie do rozwiązania dowolnej stałej zespolonej
nie zmienia ani laplasjanu, ani pochodnej normalnej.
Jeżeli więc rozwiązanie istnieje, zbiór wszystkich rozwiązań
jest dokładnie rodziną $u+\C$.
\end{proof}

\subsection{Transport: rozwiązanie przesuwa się po charakterystykach}

\begin{stwierdzenie}[Transport ze stałą prędkością]
\label{prop:transport}
Niech $c\in\R$ i $u_0\in C^1(\R)$. Funkcja
$u\colon[0,\infty)\times\R\to\R$ określona wzorem
$u(t,x)=u_0(x-ct)$ jest rozwiązaniem klasycznym
$u_t+c\,u_x=0$ z warunkiem $u(0,x)=u_0(x)$.
Jest to jedyne rozwiązanie klasy $C^1$ na
$[0,\infty)\times\R$ z tym warunkiem.
\end{stwierdzenie}
\begin{proof}
Różniczkując, dostajemy $u_t=-c\,u_0'(x-ct)$
i $u_x=u_0'(x-ct)$; równanie jest spełnione.
Jeśli $v$ jest dowolnym innym rozwiązaniem, to dla
ustalonego $\xi\in\R$ funkcja
$h(t)=v(t,\xi+ct)$ spełnia
$h'(t)=v_t(t,\xi+ct)+c v_x(t,\xi+ct)=0$.
Zatem $v(t,\xi+ct)=v(0,\xi)=u_0(\xi)$, czyli
$v(t,x)=u_0(x-ct)$.
\end{proof}

\begin{center}
\begin{tikzpicture}[>=Stealth,scale=.95]
 \draw[->] (-3.5,0)--(4.3,0) node[right] {$x$};
 \draw[->] (-3.15,-.15)--(-3.15,3.1) node[above] {$t$};
 \foreach \x in {-2.4,-1.3,-.2,.9} {
   \draw[blue!55!black,thick] (\x,0)--({\x+1.3},2.6);
 }
 \draw[densely dashed] (-3.15,2.6)--(3.7,2.6);
 \fill[blue!55!black] (-1.3,0) circle (1.8pt);
 \fill[blue!55!black] (0,2.6) circle (1.8pt);
 \node[below] at (-1.3,0) {$\xi$};
 \node[above] at (0,2.6) {$(\xi+ct,t)$};
 \node[anchor=west] at (2.6,1) {$c>0$};
\end{tikzpicture}
\par\smallskip\small
Charakterystyki równania $u_t+c\,u_x=0$ są prostymi $x=\xi+ct$;
wartość $u$ pozostaje na każdej z nich stała (rysunek dla $c>0$).
\end{center}

\begin{przyklad}[Utrata gładkości mimo gładkich danych]
\label{ex:burgers-blowup}
Rozważmy dla $t\geq0$, $x\in\R$ równanie
$u_t+u\,u_x=0$ z warunkiem $u(0,x)=-\sin x$.
Załóżmy, że rzeczywista funkcja $u$ jest rozwiązaniem klasy $C^1$ na
$[0,T)\times\R$, gdzie $T>0$. Dla $\xi\in\R$ połóżmy
$X(t)=\xi-t\sin\xi$ i $h(t)=u(t,X(t))$.
Z reguły łańcuchowej i równania Burgersa wynika
$h'(t)=-(h(t)+\sin\xi)u_x(t,X(t))$.
Funkcja $q(t)=h(t)+\sin\xi$ spełnia zatem liniowe równanie
$q'(t)=-u_x(t,X(t))q(t)$ z $q(0)=0$.
Współczynnik jest ciągły, więc po pomnożeniu przez czynnik
całkujący $\exp(\int_0^t u_x(s,X(s))\,\mathrm d\lambda_1(s))$
otrzymujemy funkcję o zerowej pochodnej i zerowej wartości
początkowej. Stąd $q(t)=0$ na całym $[0,T)$, a więc
$X'(t)=u(t,X(t))$. Otrzymujemy równanie charakterystyki
\begin{equation}\label{eq:burgers-characteristics}
 x=\xi-t\sin\xi,\qquad u(t,x)=-\sin\xi.
\end{equation}
Dla $0\leq t<1$ pochodna odwzorowania
$\xi\mapsto\xi-t\sin\xi$ wynosi $1-t\cos\xi\geq1-t>0$.
Dla ustalonego $t$ odwzorowanie $F_t(\xi)=\xi-t\sin\xi$
jest więc ściśle rosnące, a z $|F_t(\xi)-\xi|\leq t$
wynika, że jego granice przy $\xi\to\pm\infty$ są równe
odpowiednio $\pm\infty$. Z ciągłości wynika surjektywność.
Twierdzenie o funkcji uwikłanej zastosowane do równania
$F_t(\xi)-x=0$ (pochodna względem $\xi$ jest dodatnia)
daje gładką zależność odwrotności $\xi=\xi(t,x)$ od obu zmiennych,
również lokalnie w otoczeniu $t=0$.
Różniczkując $x=\xi-t\sin\xi$, dostajemy
$\xi_x=(1-t\cos\xi)^{-1}$ oraz
$\xi_t=\sin\xi/(1-t\cos\xi)$.
Dla $u=-\sin\xi$ mamy więc $u_t=-\cos\xi\,\xi_t$,
$u_x=-\cos\xi\,\xi_x$ i $u_t+uu_x=0$.
Ponadto $\xi(0,x)=x$, więc spełniony jest warunek początkowy.
W szczególności
\begin{equation}\label{eq:burgers-gradient}
 u_x(t,x)=\frac{-\cos\xi}{1-t\cos\xi}.
\end{equation}
Dla $x=0$ odpowiada $\xi=0$, a więc
$u_x(t,0)=-1/(1-t)\to-\infty$ gdy $t\uparrow1$.
Rozwiązanie pozostaje ograniczone, ale jego pochodna
przestrzenna staje się nieograniczona w skończonym czasie.
Ponieważ dla $t<1$ odwzorowanie $F_t$ jest bijekcją, każda
funkcja klasy $C^1$ spełniająca te same dane musi być równa
skonstruowanemu rozwiązaniu w zakresie $t<1$. Gdyby istniało
rozwiązanie klasy $C^1$ na $[0,T)\times\R$ dla $T>1$,
jego pochodna $u_x$ byłaby ciągła w $(1,0)$, co jest
sprzeczne z powyższą granicą.
To najprostszy mechanizm powstawania osobliwości i
fal uderzeniowych.
\end{przyklad}

\subsection{Równanie ciepła: natychmiastowe rozprzestrzenianie}

\begin{stwierdzenie}[Jądro ciepła na całej przestrzeni]
\label{prop:heat-kernel}
Niech $d\geq1$, $\kappa>0$, $u_0\in C_c^\infty(\R^d)$.
Dla $t>0$, $x\in\R^d$ zdefiniujmy
\begin{equation}\label{eq:heat-kernel}
 \Gamma_\kappa(t,x)=\frac{1}{(4\pi\kappa t)^{d/2}}
 \exp\!\left(-\frac{|x|^2}{4\kappa t}\right),\qquad
 u(t,x)=\int_{\R^d}\Gamma_\kappa(t,x-y)u_0(y)\,
 \mathrm d\lambda_d(y).
\end{equation}
Wówczas $u\in C^\infty((0,\infty)\times\R^d)$,
$u_t-\kappa\Delta_xu=0$ i
$u(t,\cdot)\to u_0$ jednostajnie, gdy $t\downarrow0$.
Jeśli dodatkowo $u_0\geq0$ i $u_0\not\equiv0$, to
$u(t,x)>0$ dla każdego $t>0$ oraz $x\in\R^d$.
\end{stwierdzenie}
\begin{proof}
Dla $t>0$ różniczkowanie wzoru na $\Gamma_\kappa$ daje
\[
 \partial_t\Gamma_\kappa=
 \left(-\frac d{2t}+\frac{|x|^2}{4\kappa t^2}\right)\Gamma_\kappa,
 \qquad
 \partial_{x_j}\Gamma_\kappa=-\frac{x_j}{2\kappa t}\Gamma_\kappa.
\]
Ponowne różniczkowanie względem $x_j$ i sumowanie po $j$
prowadzi do
$\Delta_x\Gamma_\kappa=
(-d/(2\kappa t)+|x|^2/(4\kappa^2t^2))\Gamma_\kappa$.
Stąd $\partial_t\Gamma_\kappa=\kappa\Delta_x\Gamma_\kappa$.
Niech $K$ będzie zwartym podzbiorem $(0,\infty)\times\R^d$
i niech $S=\supp u_0$. Na iloczynie odpowiednio powiększonego
zwartego otoczenia $K$ i $S$ każda ustalona pochodna
$(t,x)\mapsto\Gamma_\kappa(t,x-y)$ jest ograniczona
jednostajnie względem $(t,x,y)$: jest ciągła, a czas jest
oddzielony od zera. Iloczyn tej stałej i $|u_0(y)|$ jest
całkowalną majorantą, więc twierdzenie o różniczkowaniu całki
z parametrem stosuje się do pochodnych dowolnego rzędu.
Wynika stąd gładkość $u$ i
$u_t-\kappa\Delta_xu=\int_{\R^d}(\partial_t-\kappa\Delta_x)
\Gamma_\kappa(t,x-y)u_0(y)\,\mathrm d\lambda_d(y)=0$.

Podstawienie $z=2\sqrt{\kappa t}\,w$ ma jakobian
$(2\sqrt{\kappa t})^d$. Korzystając z twierdzenia Tonellego
dla nieujemnego iloczynu funkcji gaussowskich i wzoru
$\int_{\R}e^{-s^2}\,\mathrm d\lambda_1(s)=\sqrt\pi$, dostajemy
\[
 \int_{\R^d}\Gamma_\kappa(t,z)\,\mathrm d\lambda_d(z)
 =\pi^{-d/2}\int_{\R^d}e^{-|w|^2}\,\mathrm d\lambda_d(w)=1.
\]
Dlatego
$u(t,x)-u_0(x)=\int_{\R^d}\Gamma_\kappa(t,z)
\bigl(u_0(x-z)-u_0(x)\bigr)\,\mathrm d\lambda_d(z)$.
Funkcja $u_0$ jest jednostajnie ciągła. Ustalmy $\eps>0$
i wybierzmy $r>0$ tak, aby
$|u_0(x-z)-u_0(x)|<\eps$ dla $|z|<r$, niezależnie od $x$.
Część całki po $|z|<r$ ma moduł nie większy od $\eps$;
część po $|z|\geq r$ jest ograniczona przez
$2\|u_0\|_\infty\int_{|z|\geq r}\Gamma_\kappa(t,z)
\,\mathrm d\lambda_d(z)$, a ostatnia całka dąży do zera,
gdy $t\downarrow0$. Istotnie, po podstawieniu $z=\sqrt t\,w$
jest ona równa całce funkcji $\Gamma_\kappa(1,w)$ po
$\{|w|\geq r/\sqrt t\}$. Funkcje podcałkowe, przedłużone
przez zero na dopełnienie tego zbioru, zbiegają punktowo
do zera i są zdominowane przez całkowalną funkcję
$\Gamma_\kappa(1,w)$. Twierdzenie o zbieżności zdominowanej
daje żądaną granicę. Otrzymane oszacowanie jest niezależne
od $x$, a $\eps>0$ było dowolne, co dowodzi zbieżności
jednostajnej na $\R^d$.

Jeżeli $u_0\geq0$ i $u_0\not\equiv0$, to z ciągłości
wynika, że $u_0$ jest dodatnia na pewnym niepustym zbiorze
otwartym. Ponieważ $\Gamma_\kappa(t,x-y)>0$ dla wszystkich
$t>0$, $x,y\in\R^d$, całka we wzorze na $u(t,x)$
jest ściśle dodatnia.
\end{proof}

Jest to pierwsza istotna różnica między równaniem ciepła
a równaniem falowym: dla nieujemnych, niezerowych danych $u_0$
o zwartym nośniku, nawet bardzo małym, dla każdego $t>0$
rozwiązanie równania ciepła jest dodatnie w całej $\R^d$.
Nie oznacza to, że zaburzenie ma wszędzie porównywalną
wielkość; daleko od nośnika danych wartości są bardzo małe.

\begin{stwierdzenie}[Energetyczna jednoznaczność równania ciepła]
\label{prop:heat-uniqueness}
Niech $d\geq1$, $\Omega\subset\R^d$ będzie obszarem ograniczonym
o brzegu klasy $C^1$, $T>0$ i $\kappa>0$.
Jeśli $u,v\in C^{1,2}([0,T]\times\overline\Omega)$ są funkcjami rzeczywistymi i
spełniają $w_t-\kappa\Delta_xw=0$ dla $w=u-v$
na $(0,T)\times\Omega$, a ponadto
$u(0,\cdot)=v(0,\cdot)$ w $\Omega$ oraz
$u=v$ na $[0,T]\times\partial\Omega$, to $u=v$
na $[0,T]\times\overline\Omega$.
\end{stwierdzenie}
\begin{proof}
Niech $E(t)=\frac12\int_\Omega w(t,x)^2\,
\mathrm d\lambda_d(x)$. Funkcje $w$ i $w_t$ są ciągłe na zwartym zbiorze
$[0,T]\times\overline\Omega$, więc iloczyn $ww_t$ jest ograniczony.
Ponieważ $\Omega$ ma skończoną miarę, różniczkowanie pod znakiem
całki jest uzasadnione twierdzeniem o zbieżności zdominowanej.
Dla $0<t<T$, z równania $w_t=\kappa\Delta_xw$ i twierdzenia
Gaussa zastosowanego do pola $w\nabla_xw$ otrzymujemy
\begin{align*}
 E'(t)&=\int_\Omega w(t,x)w_t(t,x)\,\mathrm d\lambda_d(x)\\
 &=\kappa\int_{\partial\Omega}w(t,x)\partial_\nu w(t,x)\,
 \mathrm d\sigma(x)
 -\kappa\int_\Omega|\nabla_xw(t,x)|^2\,\mathrm d\lambda_d(x)\\
 &=-\kappa\int_\Omega|\nabla_xw(t,x)|^2\,\mathrm d\lambda_d(x)\leq0.
\end{align*}
Całka brzegowa znika, ponieważ $w=0$ na brzegu.
Energia jest ciągła także dla $t=0,T$. Ponieważ $E(0)=0$
i $E(t)\geq0$, jej monotoniczność daje $E(t)=0$ dla każdego
$t\in[0,T]$. Gdyby dla pewnego $t$ oraz $x\in\Omega$ zachodziło
$w(t,x)\ne0$, ciągłość dawałaby kulę zawartą w $\Omega$,
na której $w(t,\cdot)^2$ byłoby ograniczone z dołu przez dodatnią
stałą; przeczyłoby to $E(t)=0$. Zatem $w=0$ w $\Omega$
dla każdego $t$, a przez ciągłość także na $\overline\Omega$.
\end{proof}

\subsection{Równanie falowe: skończona prędkość propagacji}

\begin{twierdzenie}[Wzór d'Alemberta]
\label{thm:dalembert}
Niech $c>0$, $u_0\in C^2(\R)$ i $u_1\in C^1(\R)$.
Jedynym rozwiązaniem klasy $C^2$ na $[0,\infty)\times\R$
zagadnienia
$u_{tt}-c^2u_{xx}=0$, $u(0,x)=u_0(x)$,
$u_t(0,x)=u_1(x)$ jest
\begin{equation}\label{eq:dalembert}
 u(t,x)=\frac{u_0(x-ct)+u_0(x+ct)}{2}
 +\frac{1}{2c}\int_{x-ct}^{x+ct}u_1(s)\,\mathrm d\lambda_1(s).
\end{equation}
Jeśli $\supp u_0\cup\supp u_1\subset[-R,R]$ dla $R>0$,
to $u(t,x)=0$ dla $|x|>R+ct$.
\end{twierdzenie}
\begin{proof}
Wprowadzamy zmienne $\xi=x+ct$ i $\eta=x-ct$.
Odwzorowanie jest odwracalne, gdyż
$t=(\xi-\eta)/(2c)$ i $x=(\xi+\eta)/2$.
Dla $v(\xi,\eta)=u((\xi-\eta)/(2c),(\xi+\eta)/2)$
reguła łańcuchowa daje $\partial_t=c(\partial_\xi-\partial_\eta)$
i $\partial_x=\partial_\xi+\partial_\eta$ na funkcjach
złożonych. Ponieważ pochodne mieszane funkcji klasy $C^2$
są równe, mamy
$u_{tt}-c^2u_{xx}=-4c^2v_{\xi\eta}$ w odpowiadających sobie
punktach. Z $c>0$ wynika równoważność z równaniem
$v_{\xi\eta}=0$. Obraz półpłaszczyzny $t\geq0$ to
$\{(\xi,\eta):\xi\geq\eta\}$, nie cała płaszczyzna.
Dla każdego ustalonego $\xi$ pochodna $v_\xi(\xi,\eta)$
jest stała względem $\eta<\xi$, ponieważ $v_{\xi\eta}=0$.
Z ciągłości aż do przekątnej jest równa $a(\xi)=v_\xi(\xi,\xi)$.
Funkcja $a$ jest klasy $C^1$, więc istnieje $F\in C^2(\R)$
z $F'=a$. Dla każdego ustalonego $\eta$ funkcja
$v(\xi,\eta)-F(\xi)$ ma zerową pochodną względem $\xi>\eta$,
a z ciągłości jej wartość jest równa
$G(\eta)=v(\eta,\eta)-F(\eta)$. Funkcja $G$ jest klasy $C^2$.
Otrzymujemy zatem $v(\xi,\eta)=F(\xi)+G(\eta)$ dla $\xi\geq\eta$.
Zatem $u(t,x)=F(x+ct)+G(x-ct)$.
Warunki początkowe dają
$F(x)+G(x)=u_0(x)$ i
$cF'(x)-cG'(x)=u_1(x)$.
Różniczkując pierwszą równość i rozwiązując układ względem
$F',G'$, otrzymujemy
$F'=(u_0'+u_1/c)/2$ i
$G'=(u_0'-u_1/c)/2$.
Dla $\xi=x+ct$ i $\eta=x-ct$ korzystamy z równości
$F(\xi)+G(\eta)=F(\eta)+G(\eta)+\int_\eta^\xi F'(s)\,\mathrm d\lambda_1(s)$.
Po podstawieniu $F(\eta)+G(\eta)=u_0(\eta)$ oraz wzoru na $F'$
otrzymujemy
\[
 F(\xi)+G(\eta)=\frac{u_0(\xi)+u_0(\eta)}2
 +\frac1{2c}\int_\eta^\xi u_1(s)\,\mathrm d\lambda_1(s),
\]
czyli~\eqref{eq:dalembert}.
Dla sprawdzenia istnienia niech $H$ będzie funkcją pierwotną
$u_1$, klasy $C^2$. Prawa strona~\eqref{eq:dalembert} jest sumą
funkcji $(u_0/2+H/(2c))(x+ct)$ oraz
$(u_0/2-H/(2c))(x-ct)$, a więc jest klasy $C^2$.
Dla każdej funkcji $h\in C^2(\R)$ mamy
$\partial_t^2h(x\pm ct)=c^2h''(x\pm ct)$ i
$\partial_x^2h(x\pm ct)=h''(x\pm ct)$, więc suma spełnia
równanie falowe. Przy $t=0$ suma jest równa $u_0(x)$,
a jej pochodna czasowa jest równa $u_1(x)$.
Każde rozwiązanie musiało mieć postać~\eqref{eq:dalembert},
co dowodzi również jednoznaczności.

Jeżeli $|x|>R+ct$, to oba punkty $x-ct$, $x+ct$ leżą
po tej samej stronie przedziału $[-R,R]$, a przedział
$[x-ct,x+ct]$ nie przecina nośnika $u_1$.
Wszystkie składniki~\eqref{eq:dalembert} są więc równe zeru.
\end{proof}

\begin{center}
\begin{tikzpicture}[>=Stealth,scale=1]
 \coordinate (P) at (.7,3.0);
 \coordinate (L) at (-1.6,0);
 \coordinate (R) at (3.0,0);
 \fill[gray!14] (L)--(P)--(R)--cycle;
 \draw[->] (-4.0,0)--(4.05,0) node[right] {$x$};
 \draw[->,gray!70!black] (0,-.2)--(0,3.9) node[above,black] {$t$};
 \draw[blue!65!black,thick] (L)--(P)--(R);
 \draw[very thick] (L)--(R);
 \fill (P) circle (2pt);
 \node[above right] at (P) {$(x_0,t_0)$};
 \node[below] at (L) {$x_0-ct_0$};
 \node[below] at (R) {$x_0+ct_0$};
 \node[align=center,fill=gray!14,inner sep=2pt] at (.7,1.0) {obszar\\zależności};
\end{tikzpicture}
\par\smallskip\small
Dla równania falowego w jednym wymiarze wartość w punkcie
$(x_0,t_0)$ zależy wyłącznie od danych początkowych na
$[x_0-ct_0,x_0+ct_0]$.
\end{center}

\begin{stwierdzenie}[Lokalny bilans energii w zwężającym się przedziale]
\label{prop:pde-v31-wave-cone}
Niech $c>0$, $a,b\in\R$, $a<b$, $0<T<(b-a)/(2c)$ oraz
niech rzeczywista funkcja $u\in C^2([0,T]\times\R)$ spełnia $u_{tt}-c^2u_{xx}=0$.
Dla $0\le t\le T$ połóżmy $L(t)=a+ct$, $R(t)=b-ct$ i
\[
 E(t)=\frac12\int_{L(t)}^{R(t)}
 \bigl(u_t(t,x)^2+c^2u_x(t,x)^2\bigr)\,\mathrm d\lambda_1(x).
\]
Wtedy
\begin{equation}\label{eq:pde-v31-wave-cone}
 E'(t)=-\frac c2\bigl(u_t-cu_x\bigr)^2(t,R(t))
       -\frac c2\bigl(u_t+cu_x\bigr)^2(t,L(t))\le0.
\end{equation}
Jeśli ponadto $u(0,x)=u_t(0,x)=0$ dla $x\in[a,b]$,
to $u(t,x)=0$ dla $0\le t\le T$ i $L(t)\le x\le R(t)$.
\end{stwierdzenie}
\begin{proof}
Oznaczmy $e=(u_t^2+c^2u_x^2)/2$. Z równania falowego,
reguły iloczynu i równości pochodnych mieszanych otrzymujemy
$\partial_t e=u_tu_{tt}+c^2u_xu_{xt}
=c^2(u_tu_{xx}+u_xu_{tx})=c^2\partial_x(u_tu_x)$.
Reguła Leibniza dla całki o ruchomych końcach, z
$L'=c$ oraz $R'=-c$, daje
\[
 E'(t)=c^2[u_tu_x]_{L(t)}^{R(t)}
       -ce(t,R(t))-ce(t,L(t)).
\]
Przy prawym końcu wyrażenie $c^2u_tu_x-ce$ jest równe
$-c(u_t-cu_x)^2/2$, a przy lewym końcu wyrażenie
$-c^2u_tu_x-ce$ jest równe $-c(u_t+cu_x)^2/2$.
Otrzymujemy~\eqref{eq:pde-v31-wave-cone} dla $0<t<T$
(i z jednostronnymi pochodnymi na końcach przedziału czasu).
Jeśli dane początkowe są zerowe na $[a,b]$, to również
$u_x(0,x)=0$ na $(a,b)$, a przez ciągłość także w końcach
przedziału. Stąd $E(0)=0$.
Ponieważ $E\ge0$ i $E$ jest nierosnąca, mamy $E(t)=0$.
Ciągłość pochodnych daje $u_t=u_x=0$ w otwartym stożku
$0<t<T$, $L(t)<x<R(t)$. Dla każdego punktu tego stożka
odcinek pionowy dla $0<s\leq t$ pozostaje w otwartym stożku,
gdyż $L(s)\leq L(t)<x<R(t)\leq R(s)$. Całkowanie
$u_t(s,x)=0$ i ciągłość przy $s=0$ dają $u(t,x)=u(0,x)=0$. Ciągłość daje też równość na jego brzegu.
\end{proof}

\begin{uwaga}[Dlaczego bilans lokalny jest ważniejszy od samego wzoru]
Oszacowanie~\eqref{eq:pde-v31-wave-cone} dowodzi skończonej
prędkości propagacji bez użycia jawnego wzoru d'Alemberta.
Ta sama idea, po modyfikacji geometrii i wyrazów brzegowych,
pojawia się przy wielowymiarowym równaniu falowym i układach
hiperbolicznych; tam jawny wzór na ogół nie jest dostępny.
\end{uwaga}

\section{Poprawność w sensie Hadamarda}
\label{sec:pde-hadamard}

Wyrażenie „rozwiązać równanie” nie jest wystarczające,
jeśli nie kontrolujemy wpływu niewielkich perturbacji danych.
Ustalmy przestrzenie danych $X$ i rozwiązań $Y$, wyposażone
w określone topologie, oraz dopuszczalny zbiór danych
$D\subset X$.

\begin{definicja}[Poprawność Hadamarda]
Zagadnienie jest \emph{poprawnie postawione w sensie Hadamarda}
w parze przestrzeni $(X,Y)$ na zbiorze danych $D\subset X$,
jeśli:
\begin{enumerate}
\item dla każdego $f\in D$ istnieje rozwiązanie $u\in Y$;
\item takie rozwiązanie jest jednoznaczne w $Y$;
\item odwzorowanie rozwiązujące $S\colon D\to Y$, $S(f)=u$,
jest ciągłe względem topologii podprzestrzeni na $D$.
\end{enumerate}
W zależności od zagadnienia „dane” mogą obejmować prawą stronę,
warunki początkowe, warunki brzegowe lub ich kombinację.
\end{definicja}

Trzeci warunek jest równie istotny jak pozostałe: dane
pochodzące z pomiarów lub obliczeń numerycznych zawsze
zawierają niewielki błąd.

\begin{kontrprzyklad}[Niestabilne zagadnienie Cauchy'ego dla laplasjanu]
\label{ex:laplace-illposed}
Niech $\Omega=(0,\pi)\times(0,1)$ i rozważmy równanie
$u_{xx}+u_{yy}=0$ w $\Omega$ z danymi na dolnym brzegu
$u(x,0)=0$ i $u_y(x,0)=g(x)$.
Dla każdej liczby naturalnej $k\geq1$ określmy
\begin{equation}\label{eq:hadamard-counterexample}
 \eps_k=e^{-k/2},\qquad
 u_k(x,y)=\frac{\eps_k}{k}\sinh(ky)\sin(kx).
\end{equation}
Mamy $(u_k)_{xx}=-k^2u_k$ oraz $(u_k)_{yy}=k^2u_k$,
więc $u_k$ jest funkcją harmoniczną. Ponadto $\sinh(0)=0$ i
$\cosh(0)=1$ dają $u_k(x,0)=0$ oraz
$\partial_yu_k(x,0)=\eps_k\sin(kx)$.
Dla każdego ustalonego $r\in\N_0$ zachodzi
\[
 \|g_k\|_{C^r([0,\pi])}\leq C_r k^r e^{-k/2}\longrightarrow0,
 \qquad g_k(x)=\eps_k\sin(kx),
\]
przy stałej $C_r$ niezależnej od $k$: każda pochodna rzędu
$j\leq r$ jest ograniczona w module przez $k^j e^{-k/2}$.
Dla ustalonego $r$ wykładnik $r\log k-k/2$ dąży do $-\infty$.
A więc dane dążą do zera nawet w topologii $C^\infty([0,\pi])$,
wyznaczonej przez półnormy $C^r$ dla $r\in\N_0$.
Tymczasem na poziomie $y=3/4$ otrzymujemy
\[
 \|u_k(\cdot,3/4)\|_{L^2(0,\pi)}
 =\sqrt{\frac\pi2}\,\frac{e^{-k/2}}{k}\sinh(3k/4)
 \underset{k\to\infty}{\sim}
 \sqrt{\frac\pi8}\,\frac{e^{k/4}}{k}
 \longrightarrow\infty.
\]
W obliczeniu użyliśmy tożsamości
$\int_0^\pi\sin^2(kx)\,\mathrm d\lambda_1(x)=\pi/2$
oraz $\sinh(3k/4)=e^{3k/4}(1-e^{-3k/2})/2$.
Przyjmijmy przestrzeń danych $X=C^\infty([0,\pi])$,
a za $D\subset X$ dowolny zbiór danych zawierający $0$ i wszystkie $g_k$.
Niech $Y$ będzie przestrzenią dopuszczalnych rozwiązań klasycznej
regularności, zawierającą $0$ i wszystkie $u_k$, w której
odwzorowanie $u\mapsto u(\cdot,3/4)$ do $L^2(0,\pi)$ jest ciągłe.
Gdyby zagadnienie było poprawnie postawione w sensie Hadamarda,
jednoznaczność wymuszałaby $S(0)=0$ oraz $S(g_k)=u_k$.
Ciągłość $S$ implikowałaby wtedy zbieżność przekrojów
$u_k(\cdot,3/4)$ do zera w $L^2(0,\pi)$, sprzeczną z obliczeniem.
Tak więc poprawność w sensie Hadamarda nie zachodzi w tych
przestrzeniach; jeśli rozwiązania są jednoznaczne, zawodzi
ciągła zależność od danych. Prosta $y=0$ jest niecharakterystyczna
dla operatora Laplace'a przedłużonego na otoczenie brzegu,
gdyż jej kowektor normalny $(0,1)$ nie zeruje symbolu.
\end{kontrprzyklad}

\begin{uwaga}[Dlaczego to nie przeczy zasadzie maksimum?]
We wniosku~\ref{cor:dirichlet-stability} podawaliśmy
\emph{wartości funkcji na całym brzegu}.
W kontrprzykładzie~\ref{ex:laplace-illposed} zadaliśmy
\emph{wartość i pochodną normalną tylko na części brzegu}.
To dwa zupełnie różne zagadnienia. W szczególności
zasada maksimum nie oszacowuje rozwiązania w prostokącie
przez dane Cauchy'ego wyłącznie na dolnej krawędzi.
\end{uwaga}

\begin{kontrprzyklad}[Równanie ciepła odwrócone w czasie]
\label{ex:backward-heat}
Niech $\kappa>0$, $T>0$, $\T=\R/(2\pi\mathbb Z)$.
Rozważmy na $[0,T]\times\T$ równanie
$u_t=-\kappa u_{xx}$. Dla każdej liczby całkowitej $k\geq1$
funkcja
\[
 u_k(t,x)=\exp\!\bigl(\kappa k^2(t-T/2)\bigr)\sin(kx)
\]
jest dobrze określona na torusie, ponieważ $k$ jest liczbą całkowitą.
Mamy $(u_k)_t=\kappa k^2u_k$ i $(u_k)_{xx}=-k^2u_k$, zatem
jest to gładkie rozwiązanie. Jej dane początkowe
$u_k(0,x)=e^{-\kappa k^2T/2}\sin(kx)$ dążą do zera
w $C^\infty(\T)$: pochodna rzędu $j$ ma normę supremum
nie większą niż $k^je^{-\kappa k^2T/2}$, która dla ustalonego
$j$ dąży do zera. Na torusie przyjmujemy miarę będącą obrazem
miary Lebesgue'a z $[0,2\pi)$ przez odwzorowanie ilorazowe;
jej masa wynosi $2\pi$. Ponieważ
$\int_0^{2\pi}\sin^2(kx)\,\mathrm d\lambda_1(x)=\pi$,
mamy $\|u_k(T,\cdot)\|_{L^2(\T)}
=\sqrt\pi\,e^{\kappa k^2T/2}\to\infty$.
Dla przestrzeni danych z topologią $C^\infty(\T)$ i przestrzeni
rozwiązań zawierającej $0,u_k$ oraz kontrolującej normę
$L^2(\T)$ w chwili $T$, jednoznaczność i ciągła zależność od
danych nie mogą zachodzić jednocześnie, zgodnie z argumentem
z poprzedniego kontrprzykładu. Zatem problem ewolucji wstecznej
nie jest poprawnie postawiony w sensie Hadamarda;
formalna zmiana znaku przy laplasjanie odwraca mechanizm
wygładzania w gwałtowne wzmacnianie wysokich częstości.
\end{kontrprzyklad}

\section{Skalowanie i związek typu równania z geometrią}
\label{sec:pde-scaling}

Dla równań ewolucyjnych czas i przestrzeń nie zawsze
skalują się w ten sam sposób. Współczynniki $c$ i $\kappa$
mają odpowiednio wymiary długość/czas oraz
długość$^2$/czas. To sugeruje dwa odmienne skalowania.

\begin{stwierdzenie}[Skalowania równań ciepła i falowego]
Niech $\lambda,\kappa,c>0$, $d\geq1$ i niech
$Q\subset\R\times\R^d$ będzie zbiorem otwartym.
\begin{enumerate}
\item Jeśli $u$ jest rozwiązaniem klasycznym
$u_t-\kappa\Delta_xu=0$ na $Q$, to
$u_\lambda(t,x)=u(\lambda^2t,\lambda x)$
rozwiązuje to samo równanie na
$Q_\lambda=\{(t,x):(\lambda^2t,\lambda x)\in Q\}$.
\item Jeśli $u$ jest rozwiązaniem klasycznym
$u_{tt}-c^2\Delta_xu=0$ na $Q$, to
$v_\lambda(t,x)=u(\lambda t,\lambda x)$
rozwiązuje to samo równanie na
$\{(t,x):(\lambda t,\lambda x)\in Q\}$.
\end{enumerate}
\end{stwierdzenie}
\begin{proof}
Dziedziny funkcji przeskalowanych są otwarte jako przeciwobrazy
$Q$ przez ciągłe odwzorowania liniowe. Ponieważ odwzorowania
te osobno skalują czas i zmienne przestrzenne, zachowana
jest wymagana regularność klasyczna. Reguła łańcuchowa daje
$\partial_tu_\lambda=\lambda^2u_t(\lambda^2t,\lambda x)$
i $\Delta_xu_\lambda=\lambda^2\Delta_xu(\lambda^2t,\lambda x)$.
Stąd
$(\partial_t-\kappa\Delta_x)u_\lambda(t,x)
=\lambda^2(u_t-\kappa\Delta_xu)(\lambda^2t,\lambda x)=0$.
Podobnie $\partial_t^2v_\lambda=\lambda^2u_{tt}(\lambda t,\lambda x)$
oraz $\Delta_xv_\lambda=\lambda^2\Delta_xu(\lambda t,\lambda x)$,
więc $(\partial_t^2-c^2\Delta_x)v_\lambda(t,x)
=\lambda^2(u_{tt}-c^2\Delta_xu)(\lambda t,\lambda x)=0$.
\end{proof}

Przy skali przestrzennej $\ell>0$ podstawienie $x=\ell z$,
$t=\tau s$ do równania ciepła daje współczynnik
$\kappa\tau/\ell^2$ przy laplasjanie w zmiennej $z$.
Wybór $\tau=\ell^2/\kappa$ sprowadza go do jedynki.
Dla równania falowego analogiczny współczynnik wynosi
$c^2\tau^2/\ell^2$, więc naturalną skalą czasu jest $\ell/c$.
To różnica strukturalna, nie tylko rachunkowa.

\section{Co będzie potrzebne dalej?}
\label{sec:pde-outlook}

Podsumujmy rozróżnienia, które pozostaną ważne przez cały skrypt:
\begin{enumerate}
\item \emph{Równanie a zagadnienie.} Dwa takie same operatory
mogą prowadzić do różnych problemów przy innych danych
początkowych lub brzegowych.
\item \emph{Istnienie, jednoznaczność, stabilność.} Każde z tych
zagadnień wymaga osobnego rozumowania. Udowodnienie
jednoznaczności nie implikuje istnienia.
\item \emph{Typ operatora.} Eliptyczność, hiperboliczność
i paraboliczność opisują zasadniczo różne mechanizmy.
\item \emph{Klasa rozwiązań.} Rozwiązania klasyczne są tylko
początkiem. Naturalnymi obiektami bywają funkcje Sobolewa,
dystrybucje oraz rozwiązania słabe i lepkościowe.
\item \emph{Rola geometrii.} Symbol główny, charakterystyki,
struktura brzegu i skalowanie decydują o propagacji
i regularności rozwiązań.
\end{enumerate}

W następnym rozdziale rozwiniemy metodę charakterystyk dla
równań pierwszego rzędu: od liniowego transportu przez
quasiliniowe zagadnienia Cauchy'ego po tworzenie się
fal uderzeniowych.

% Rozdział 2.
